\uuid{YyZh}
\exo7id{3313}
\titre{exo7 3313}
\auteur{quercia}
\organisation{exo7}
\datecreate{2010-03-09}
\isIndication{false}
\isCorrection{false}
\chapitre{Application linéaire}
\sousChapitre{Image et noyau, théorème du rang}
\cpge{mpsi-mp2i.applications-lineaires.e, pcsi.applications-lineaires.f, mpsi-mp2i.espaces-vectoriels.d, pcsi.espaces-vectoriels.d, ptsi.espaces-vectoriels.d}
\module{Algèbre}
\niveau{L1}
\difficulte{}

\contenu{
\texte{
Soit $E$ un $ K$-ev et $f : E \to { K}$ une forme linéaire non identiquement
nulle. On note $H = \mathrm{Ker} f$.
}
\begin{enumerate}
    \item \question{Montrer que $\Im f =  K$.}
    \item \question{Soit $\vec u \in E \setminus H$ et $F = \text{vect}(\vec u)$.
    Montrer que $F \oplus H = E$.}
\end{enumerate}
}
